% src: PREREGISTRATION_ROUND2.md § 1. Purpose and relation to round 1; § 2. What had been seen when this was written; DEVIATIONS.md DEV-9.
\subsection{Round 2: registered extension}
\label{sec:r2design}
Round~2 tests the scope of the released-checkpoint findings through detector changes, retraining,
a silence-removal intervention and additional channel conditions. The registration and analysis code were
frozen together at \texttt{round2-freeze}, before any round-2 development, evaluation or target score existed.
The design was informed by the completed round-1 analyses and published detector results; pre-freeze
score-pipeline smoke tests used LA19 training clips only. No audio for the additional stress conditions
had been generated. The round-1 roster, hypotheses and results remain fixed. The extension has its own
confirmatory families and descriptive reading rule; it does not add models to any round-1 hypothesis count.

\subsubsection{Detectors and evaluation surfaces}
% src: PREREGISTRATION_ROUND2.md § 3. Models; § 5. Data and targets.
The extension adds the ASVspoof~5 baseline AASIST, the logistic-regression detector of
\citet{pascu24_calibrated}, and XLSR-Mamba \citep{xiao25_xlsrmamba}. The first uses the released ASVspoof~5
checkpoint, trained on the training partition without the post-processing applied to the development and
evaluation partitions \citep{wang25_asvspoof5_db}. Pascu's detector uses frozen XLS-R~2B features
\citep{babu21_xlsr}, mean-pools the last hidden state and fits a logistic-regression head. We refit that head
on LA19 train alone, reserving LA19 dev for threshold setting; a separate train-plus-dev fit is used only
for the reproduction gate. XLSR-Mamba uses the released deepfake-detection checkpoint trained on LA19
with RawBoost augmentation \citep{tak22_rawboost}. Table~\ref{tab:r2matrix} gives the registered scoring
matrix and source of each model's thresholds.

% src: PREREGISTRATION_ROUND2.md § 3. Models; § 5. Data and targets; § 9. Exploratory analyses (E3, E6, E11).
\begin{table}[!t]
\caption{Registered round-2 model--surface matrix. LA19 eval and ITW are common targets; additional
targets are restricted as shown. ``Stress'' denotes the common clean manifest and channel variants.
The gate-only fit contributes no threshold-transfer result.}
\label{tab:r2matrix}
\centering
\footnotesize
\setlength{\tabcolsep}{4pt}
\begin{tabular}{@{}p{0.22\textwidth}p{0.20\textwidth}p{0.23\textwidth}p{0.27\textwidth}@{}}
\toprule
Model & Threshold source & Common targets & Additional targets \\
\midrule
ASVspoof~5 AASIST & ASVspoof~5 dev & LA19 eval, ITW, stress & Verified DF21, MLAAD-5, ASVspoof~5 eval \\
Pascu LR, train-only fit & LA19 dev & LA19 eval, ITW, stress & Verified DF21, MLAAD-5 \\
XLSR-Mamba & LA19 dev & LA19 eval, ITW, stress & Verified DF21, MLAAD-5 \\
AASIST-L, untrimmed runs & Each run's LA19 dev & LA19 eval, ITW, stress & None \\
AASIST-L, trimmed runs & Each run's trimmed LA19 dev & LA19 eval, ITW, stress & None \\
Last-epoch copies of both arms & Corresponding source dev & LA19 eval, ITW & None; selection-free analysis only \\
Pascu LR, train-plus-dev fit & None & LA19 eval for the gate & None \\
\bottomrule
\end{tabular}
\end{table}

% src: PREREGISTRATION_ROUND2.md § 3. Models; § 6. Scoring, gates and thresholds.
All scoring uses desktop fp32 arithmetic with TF32 disabled and the shared score convention of
Section~\ref{sec:protocol}. The ASVspoof~5 baseline and the AASIST-L runs use 64,600 samples;
XLSR-Mamba uses 66,800 samples on every surface, including ITW. Evaluation retains the start of long
clips and repeat-tiles short clips according to the author implementations. Pascu's detector uses
the whole utterance; its score is the bona fide log-odds, obtained by reversing the sign of the fitted
classifier's fake-class decision function. These input and score rules are fixed across targets.

% src: PREREGISTRATION_ROUND2.md § 6. Scoring, gates and thresholds; § 10. Failure, missing data and stopping rules.
The reproduction gates check pipeline fidelity, rather than transfer performance. The ASVspoof~5
baseline must reproduce its published development EER within 15\% relative error. XLSR-Mamba is
compared with the authors' released deepfake scores on the same verified DF21 trials and at the same
input length, using the same relative tolerance. Pascu's separate train-plus-dev fit must reproduce
the published LA19-eval EER within that tolerance; it is reported only as a gate, not as an audited
detector. The training runs have no reproduction gate. A gate failure remains visible in exploratory
reporting and counts against support of the detector hypothesis; it does not trigger checkpoint or
model substitution. Every operational model receives all five registered policies from its own source
dev: \(\tauE\), \(\tauF{1}\), \(\tauF{5}\), \(\tauA{1}\) and \(\tauA{5}\). Those thresholds are frozen
once written, with no target-dependent tuning.

\subsubsection{Retraining and the silence intervention}
% src: PREREGISTRATION_ROUND2.md § 3. Models; § 4. Training runs.
We retrain the author's AASIST-L recipe \citep{jung22_aasist} with seeds 1--5 in each of two arms:
the original audio and audio trimmed by the fixed energy-VAD rule. An additional untrimmed run at the
author's default seed, 1234, checks the pipeline and enters the exploratory seed comparison only.
The schedule is held fixed: 100 epochs, batches of 24, Adam with an initial learning rate of
\(10^{-4}\), per-step cosine annealing to \(5\times10^{-6}\), weight decay \(10^{-4}\), and the recipe's
class-weighted cross-entropy. Evaluation data are excluded from training and checkpoint selection;
the author's evaluation-based best-checkpoint step is omitted.

% src: PREREGISTRATION_ROUND2.md § 4. Training runs.
The primary checkpoint is the recipe's final stochastic weight average (\texttt{swa.pth}). The
selection rule is preserved, including its initial incumbent dev EER of 1 on a percentage scale:
only epochs selected by the dev-improvement rule at an EER no greater than 1\% can enter the average.
If none is selected, the primary checkpoint is the last epoch. We report the number of snapshots
contributing to each average. The initial incumbent is a selection state, not an observed minimum
dev EER; the latter must be taken from the epoch log. Last-epoch weights are saved separately for
the selection-free comparison, and are identical to the primary weights when the average has no
snapshots. This comparison tests whether any difference between arms depends on dev-based
checkpoint selection.

% src: PREREGISTRATION_ROUND2.md § 4. Training runs (trimmed arm); § 10. Failure, missing data and stopping rules.
The VAD computes frame energy from float32 samples with 25-ms frames and a 10-ms hop. Frames within
35~dB of the clip's loudest frame are active. Both ends are cut with no added margin; a clip is left
untrimmed if the retained region would be shorter than 0.25~s. In the trimmed arm, this operation
precedes the recipe's random training crop or evaluation padding on train and dev, and precedes
padding for every scored input. Thus the intervention changes the full training-and-scoring pipeline
while keeping the detector recipe fixed. The comparison concerns this recipe and this explicit
silence rule. A pre-training check refuses any trimmed training clip of exactly 64,600 samples,
which would otherwise hit the recipe's undefined random-crop boundary.

\subsubsection{Channel intervention}
% src: PREREGISTRATION_ROUND2.md § 5. Data and targets.
The channel experiment uses the frozen, class-balanced 1,000-clip LA19-eval manifest from H5 and a
byte-identical clean copy. It adds Opus at 6 and 24~kbps, AMR-NB at 4.75~kbps, MP3 at 32~kbps,
MUSAN noise \citep{snyder15_musan} at 20, 10 and 5~dB SNR, and a mock-up-phone impulse response
from AIR within RIRS\_NOISES \citep{jeub09_air,ko17_rirs}. Codec outputs are decoded to 16-kHz PCM;
AMR-NB coding uses 8~kHz internally. Noise-file selection and offset are deterministic functions of
the utterance identifier, and SNR is defined over the whole clip. The impulse response is chosen in
the same deterministic way, its channels are averaged to mono and the response is peak-normalised
before convolution. The output is cut to the input length and peak-limited if needed to avoid clipping.
A response at an incompatible sample rate is refused,
and the chosen files and hashes are retained as provenance. The existing EnCodec condition is kept
for exploratory comparison. Every round-1 model, round-2 detector and primary training-run checkpoint
is scored on the stress set; last-epoch copies are excluded. These are channel perturbations of the
same labelled clips, so bona fide labels are retained.

\subsubsection{Confirmatory tests and their reading}
% src: PREREGISTRATION_ROUND2.md § 5. Data and targets; § 7. Confirmatory hypotheses (common rules).
For target-level tests, speakers are resampled with replacement and all their clips of both classes
are retained. ITW uses its supplied speaker labels and LA19 eval uses its corpus speakers. Each
bootstrap test uses \(B=10{,}000\), seed 0, and the registered one-sided probability
\(p=(1+\#\{T_b\leq0\})/(B+1)\), with the tested statistic centred at the null boundary.
Holm's procedure controls \(\alpha=0.05\) within each specified family \citep{holm79_holm}.
The following definitions fix the support rules and interpretation before outcomes are read.

\begin{description}
% src: PREREGISTRATION_ROUND2.md § 7. Confirmatory hypotheses (R2-H1).
\item[R2-H1: policy and error direction.]
For each new detector on ITW, test \(\frr-\far>0\) at \(\tauF{1}\), and
\(\far-\frr>0\) at \(\tauA{5}\). Holm correction covers all six tests. A detector supports the
hypothesis only when both nulls are rejected; overall support requires at least two of the three
detectors. A failed gate or an unscorable detector counts as not supporting it. Failure of the
overall rule means that transferred error direction depends on the detector as well as the policy.

% src: PREREGISTRATION_ROUND2.md § 7. Confirmatory hypotheses (R2-H2).
\item[R2-H2: transfer gap across retraining.]
For each untrimmed AASIST-L run, test whether LA19-eval TTG at \(\tauE\) exceeds 1~pp,
recomputing the oracle HTER in every speaker-cluster replicate. Holm correction covers the five
runs, and support requires at least four rejections. A missing run counts as not rejecting.
Support would extend the released-checkpoint transfer-gap finding to repeated training of this
recipe at the registered margin.

% src: PREREGISTRATION_ROUND2.md § 7. Confirmatory hypotheses (R2-H3).
\item[R2-H3: training without the silence cue.]
Compare the trimmed and untrimmed run means using exact one-sided permutation tests over the 252
reassignments under the full planned five-per-arm design; with eligible missing runs, enumerate the assignments
of the remaining run counts. Higher trimmed-arm LA19-dev EER is the manipulation check,
outside the correction family. The family contains lower ITW FRR at \(\tauE\) and lower ITW EER in
the trimmed arm. Rejection for FRR alone indicates an operating-point effect; rejection for both FRR and EER indicates
an additional effect on discrimination. If neither null is rejected, the registered reading
is that the intervention does not explain the recipe's ITW behaviour. With fewer than three completed runs in either arm, the hypothesis is
not evaluable; the registered analysis uses the remaining runs when that minimum is met.

% src: PREREGISTRATION_ROUND2.md § 7. Confirmatory hypotheses (R2-H4).
\item[R2-H4: additional channels.]
For each new condition and each primary round-1 model, test
\(\Delta\frr-\Delta\far>0\) at the frozen round-1 \(\tauE\), relative to the clean copy of the
same clips. Clean and stressed scores remain paired in speaker resampling. Holm correction covers
the four models within each condition, and a condition supports the hypothesis if at least three
models reject. The component tests \(\Delta\frr>0\) and \(\Delta\far<0\) are corrected and
reported in the same way. Reporting is by condition, including the number of conditions meeting
the rule, rather than a pooled claim across channels.
\end{description}

% src: PREREGISTRATION_ROUND2.md § 8. Reading rule R2-RL (fixed in advance, descriptive).
The descriptive reading rule applies to each new detector on ITW at \(\tauE\) in a fixed order.
An ITW EER of at least 20\% is \emph{discrimination-limited}. Otherwise, a detector is
\emph{transfer-limited} if its TTG/HTER is at least 0.5 or its source-dev-equivalent threshold
interval spans at least 10~pp of ITW FRR. It is \emph{well-transferred} if the share is below 0.2
and the span below 5~pp, and \emph{intermediate} otherwise. These categories describe whether the
round-1 pattern carries over; they are not additional hypothesis tests.

\subsubsection{Exploratory analyses}
% src: PREREGISTRATION_ROUND2.md § 9. Exploratory analyses (E1--E6).
The exploratory programme reports effect sizes and intervals without confirmatory language.
E1 describes FRR, FAR, HTER, oracle HTER, TTG, transfer share, dev-equivalent intervals and their
target endpoints, and the rank-based miss and band on every registered model--surface--policy cell,
using speaker-cluster 95\% intervals with \(B=1{,}000\). E2 places the released AASIST-L values
within the spread of the retrained arms, including the pipeline-check run. E3 compares last-epoch
copies without checkpoint selection, and E4 reports paired channel effects for every scored model
and run, including EnCodec. E5 reports the new detectors on MLAAD-5 by language and states its
sample-rate confound with every numerical result: native sample rate is associated with class,
as quantified by the round-1 census in Table~\ref{tab:r1mlaad}. The already-seen round-1 MLAAD results remain descriptive. E6 reports
the ASVspoof~5 baseline by evaluation condition, with codec-free C00 primary, and characterises
development and evaluation silence with the same VAD.

% src: PREREGISTRATION_ROUND2.md § 9. Exploratory analyses (E7, E8).
E7 evaluates development-designed accept/review/reject policies. The primary review budget is
10\% of pooled dev traffic, centred in rank at \(\tauE\), with edges at
\(F_{\mathrm{dev}}(\tauE)\pm0.05\). Each target contributes its pooled and class-specific review
shares, bona fide auto-rejection and spoof auto-acceptance. Comparators use the same construction
at the labelled target oracle and a label-free variant re-anchored on the target's pooled scores.
Secondary budgets are 5\% and 20\%. A separate error-budget rule rejects below \(\tauF{1}\) and
accepts from \(\tauA{1}\); if those bounds cross or coincide, it collapses to \(\tauE\), and the
collapse is reported. E8 gives the APCER/FAR and BPCER/FRR reading at every policy, including
BPCER at source-fixed APCER budgets of 1\% and 5\% and at the corresponding target oracle.

% src: PREREGISTRATION_ROUND2.md § 9. Exploratory analyses (E9--E11); DEVIATIONS.md DEV-9 amendment 2.
E9 compares actual and minimum normalised cost after an unpenalised logistic calibration fitted
on source dev alone, with prior correction and Bayes thresholds over logit priors from \(-4\) to
\(4\) in steps of 0.25. The logged amendment replaces the registered fixed-step fitting routine
with the maximum-likelihood logistic fit after a unit test exposed scale-dependent underfitting;
the prior sweep and cost comparison are unchanged. E10 reports the screening and forensic
break-even windows at \(\tauE\), using the round-1 cost models. E11 contains the train-plus-dev
Pascu fit solely as the reproduction gate.

\subsubsection{Failures and amendments}
% src: PREREGISTRATION_ROUND2.md § 10. Failure, missing data and stopping rules; DEVIATIONS.md DEV-9 amendment 1.
The failure rules prohibit replacement detectors, additional seeds and repetition after training
divergence, NaN loss or a training crash. Missing runs count against the registered retraining
support rule; the silence test uses only the completed runs subject to its evaluability minimum.
A logged resource amendment permits a run that fails before completing its first epoch, before
any dev score exists, to restart once from scratch with the same seed, code and data. Scheduling
was adjusted using observed memory requirements, with failed directories retained. This exception
does not permit restarting a run after training outcomes have been observed. A logistic-regression
convergence warning is reported and the fitted head is retained.

% src: PREREGISTRATION_ROUND2.md § 10. Failure, missing data and stopping rules; DEVIATIONS.md DEV-9 amendment 3.
Loaders refuse unknown labels or codecs, incorrect row counts and missing files. Every such fix
is logged before rerunning the affected job. The FLAC amendment requires native or fallback
integer decoding to match the file's STREAMINFO MD5 before explicit float32 conversion. It was
logged after a decoder sign error was found, before any affected round-2 DF21 score was analysed;
the affected scores were quarantined for re-scoring. This preserves the registered verified
surface and prevents an environment-specific decoder from silently changing its audio.

% src: DEVIATIONS.md DEV-9 amendment 4.
A separate amendment permits one unchanged rerun of Pascu's whole-utterance MLAAD-5 scoring on a
GPU free of a training run, with an expandable memory-allocator setting. It does not shorten
utterances, change the head or alter arithmetic. A second failure is reported as missing MLAAD-5
data; that surface contributes only to the exploratory language analysis. Each amendment records
what had been seen when it was made. These resource and fidelity amendments are reported alongside
the frozen design, rather than used to revise its tests or support thresholds.
